\documentclass[11pt, a4paper]{report}

\usepackage[utf8]{inputenc}
\usepackage[T1]{fontenc}
\usepackage[french]{babel}

% Appel des packages dans le dossier preambule
\def\packagepath{../../../../../preambule}   % path du package princial

\usepackage{\packagepath/page_layout} % <--- Nouveau
\usepackage{\packagepath/config_info}
\usepackage{\packagepath/cadres_cours}
\usepackage{\packagepath/zones_reponse}
\usepackage{\packagepath/config_ds}
\usepackage{\packagepath/config_maths}

% --- PILOTAGE ---
%\def\visibleornot{invisible}    % visible or invisible
\ifdefined\visibleornot\else\def\visibleornot{visible}\fi


\def\documentpath{./Sources_Latex}
\graphicspath{{./Images/}}


\def\ptsexoA{0}



\begin{document}

\def\level{Premiere}  % Classe à droite
\def\course{NSI}
\def\chaptername{Les algorithmes de recherche} % Titre au centre
\def\docname{AP06~--~TD 1}        % Identifiant à gauche

\pagestyle{DS_LP}

\NomPrenom{}

\begin{exercice_cours}[Complexité des algorithmes]{}
    
    \begin{enumerate}
        \item Si un algorithme a une complexité de $O(n)$, que se passe-t-il pour le temps d'exécution si la taille de la liste est multipliée par 10? \hfill
        \reponseLigne[10cm]{\texttt{Le temps d'exécution est multiplié par 10.}}
        \item Si un algorithme a une complexité de $O(n^2)$, que se passe-t-il pour le temps d'exécution si la taille de la liste est multipliée par 10? \hfill
        \reponseLigne[10cm]{\texttt{Le temps d'exécution est multiplié par 100.}}
        \item Si un algorithme a une complexité de $O(\log n)$, que se passe-t-il pour le temps d'exécution si la taille de la liste est multipliée par 64? \hfill
        \reponseLigne[10cm]{\texttt{Le temps d'exécution augmente de 6.}}
    \end{enumerate}   
\end{exercice_cours}

\begin{exercice_cours}[On considère une liste triée \texttt{L}.]
    
    \begin{enumerate}
        \item     Ecrire un algorithme de recherche dichotomique pour trouver la position d'un élément \verb|x| dans la liste. Si l'élément n'est pas présent, l'algorithme doit retourner \verb|-1|.


\item Indiquer quel est le variant de boucle de cet algorithme. \hfill\reponseLigne[3cm]{\texttt{right - left}}
    \item Indiquer quel est l'invariant de boucle de cet algorithme. \hfill
    
    \reponseLigne[0.8\linewidth]{\texttt{L[left:right] est triée et contient x si x est présent}}
    \item Quelle est la complexité de cet algorithme? \hfill\reponseLigne[3cm]{$O(\log n)$}  

    \item Combien de valeurs sont examinees lors d’un appel à 
    
    \verb|recherche_dichotomique|([1, 3, 4, 6, 9, 12, 14, 15],9)| et quelle est la valeur renvoyée ?
    
            \reponseLigne[0.8\linewidth]{}

    \item Même question pour \verb|recherche_dichotomique([1, 3, 4, 6, 9, 12, 14, 15], 7)|.

            \reponseLigne[0.8\linewidth]{}

\end{enumerate}

\end{exercice_cours}


\begin{minipage}{0.48\linewidth}
        \begin{exercice_cours}[Débogage]
            \begin{enumerate}
                \item Le code suivant pour la recherche naïve contient une erreur classique. Laquelle ?
                \item Corriger le code pour qu'il fonctionne correctement.
            \end{enumerate}
    \end{exercice_cours}

    \end{minipage}\hfill
    \begin{minipage}{0.48\linewidth}
    \begin{CodePythontexAlt}[TaillePolice=\large, Largeur=1\linewidth,Centre,Lignes=true,PremLigne=1]{}
def recherche(elt, L):
    for x in L:
        if x == elt:
            return True
        else:
            return False
    \end{CodePythontexAlt}
    \end{minipage}



    \begin{minipage}{0.48\linewidth}
        \begin{exercice_cours}[Débogage]
         Le code suivant présente une erreur, l'identifier et la corriger.
    \end{exercice_cours}

    \end{minipage}\hfill
    \begin{minipage}{0.48\linewidth}
    \begin{CodePythontexAlt}[TaillePolice=\large, Largeur=1\linewidth,Centre,Lignes=true,PremLigne=1]{}
i = 10
while i > 0:
    print(i)
    \end{CodePythontexAlt}
    \end{minipage}

    

\end{document}